% !TEX root = ../main.tex
\chapter{Sparse Identification of Nonlinear Dynamics}\label{ch:sindy}
\begin{paracol}{2}

\begin{sources}
\begin{tabularx}{\linewidth}{@{}V l L@{}}
\toprule
\textbf{Video} & \textbf{Length} & \textbf{Content}\\
\midrule
\seg{5}{1}{V5.1 Sparse Identification of Nonlinear Dynamics (SINDy)} & 26:44 & The SINDy regression on Lorenz; noisy data and TV derivatives; the cylinder wake and the mean-field model (history from Ruelle--Takens to Noack); kicking off the attractor; parametrized, time-dependent and forced models; SINDy in delay coordinates; SINDy with control; open questions.\\
\seg{5}{2}{V5.2 PDE FIND} & 10:36 & (Speaker: Samuel Rudy; uploaded by N.~Kutz.) Discovering PDEs $u_t=N(u,u_x,\dots)$ by sparse regression; derivatives; sequential threshold ridge regression; Navier--Stokes from 1\% of the data; the soliton ambiguity (advection vs KdV); a gallery of PDEs.\\
\bottomrule
\end{tabularx}

\smallskip
\textbf{Transcripts} \file{materials/transcripts/C05-V0*.txt}.
\textbf{Book} Sec.~7.3 (SINDy); Sec.~3.5 (sparse regression); Sec.~4.4--4.7 (regression, model selection).
\textbf{Papers} S.~L.~Brunton, J.~L.~Proctor, J.~N.~Kutz, PNAS 113 (2016) 3932--3937; S.~H.~Rudy,
S.~L.~Brunton, J.~L.~Proctor, J.~N.~Kutz, Sci.\ Adv.\ 3 (2017) e1602614.
\textbf{Code} \file{code/ch05_sindy.py}, \file{ch05_pdefind.py}.
\end{sources}

\section*{Motivation}
DMD (\cref{ch:dmd}) fits a \emph{linear} model to data. SINDy goes after the second primary challenge of
\cref{ch:overview}, unknown dynamics, head on: it discovers the nonlinear right-hand side $\vb f$ itself,
as a short formula, from measurements \ts{5}{1}{0:04}. If we had $\vb f$ we could integrate trajectories,
compute coherent structures and FTLEs, quantify uncertainty \ts{5}{1}{0:45}\ts{5}{1}{1:06}; in neuroscience or
climate science we do not have it, so we want to discover it from data \ts{5}{1}{1:27}. The key assumption is
\emph{parsimony}: in the right coordinates, $\vb f$ has only a few terms.

% ------------------------------------------------------------------
\section{The SINDy algorithm}\label{sec:si-algorithm}

\paragraph{Sparsity in function space.} In the Lorenz equations \eqref{eq:si-lorenz} the $\dot x$ equation
contains only $x$ and $y$; $\dot y$ contains $x$, $xz$ and $y$; $\dot z$ contains $z$ and $xy$
\ts{5}{1}{1:49}\ts{5}{1}{2:09}. Out of all possible right-hand sides this one is very special: it is
\emph{sparse} in the space of candidate functions \ts{5}{1}{2:30}\ts{5}{1}{3:11}.

\begin{definition}[New concepts: library, SINDy]\label{def:sindy}
From $m$ samples of the state and its derivative, build
\[
  X=\begin{bmatrix}\vb x^{\mathsf T}(t_1)\\\vdots\\\vb x^{\mathsf T}(t_m)\end{bmatrix}\in\R^{m\times n},\qquad
  \dot X\in\R^{m\times n},
\]
and a \textbf{library} of candidate functions evaluated on the data,
$\Theta(X)=\begin{bmatrix}\vb 1 & X & X^{P_2} & X^{P_3}&\cdots\end{bmatrix}$, where $X^{P_2}$ holds all
quadratic monomials ($x^2,xy,xz,y^2,\dots$), and so on \ts{5}{1}{3:32}. \textbf{SINDy} (Sparse
Identification of Nonlinear DYnamics) solves
\begin{equation}\label{eq:sindy}
  \dot X\approx\Theta(X)\,\Xi,\qquad \Xi=[\boldsymbol\xi_1\cdots\boldsymbol\xi_n]\ \text{sparse},
\end{equation}
one sparse regression per state variable: $\dot{\vb x}_k\approx\Theta(X)\boldsymbol\xi_k$ \ts{5}{1}{3:52}\ts{5}{1}{4:12}.
The identified model is $\dot x_k=\boldsymbol\theta(\vb x^{\mathsf T})\boldsymbol\xi_k$, with
$\boldsymbol\theta$ the row of library functions.
\end{definition}

\paragraph{Why sparse regression.} Plain least squares would use ``a little bit of all'' the columns
\ts{5}{1}{5:34}; an $\ell_1$-penalized (or thresholded) regression finds a good fit with as few terms as
possible \ts{5}{1}{5:56}. Options: the LASSO of Tibshirani, or the \emph{sequentially thresholded least
squares} of the paper \ts{5}{1}{4:53}.

\begin{definition}[New concepts: LASSO, sequentially thresholded least squares]
\begin{enumerate}[(a)]
  \item \textbf{LASSO} (Least Absolute Shrinkage and Selection Operator, Tibshirani 1996):
        $\min_{\boldsymbol\xi}\lVert\dot{\vb x}-\Theta\boldsymbol\xi\rVert_2^2+\alpha\lVert\boldsymbol\xi\rVert_1$.
  \item \textbf{STLSQ}: solve least squares; set to zero every coefficient with
        $|\xi_i|<\lambda$; solve least squares again on the surviving columns; repeat until the support
        stops changing. The threshold $\lambda$ is the sparsity knob.
\end{enumerate}
\end{definition}

\codefile{SINDy on the Lorenz system}{../code/ch05_sindy.py}

The script, with derivatives computed from the data by finite differences and a library of all 56
monomials up to order 5, prints
\[
  \dot x=-10.000x+10.000y,\quad \dot y=27.998x-1.000y-1.000xz,\quad \dot z=-2.667z+1.000xy ,
\]
the Lorenz system with its parameters; least squares would have kept 133 of the 168 coefficients. As in
the video: same structure, virtually the same $\sigma,\rho,\beta$, same dynamics on the attractor
\ts{5}{1}{6:16}.

\begin{intuition}[Occam's razor as an optimization]
The threshold implements a Pareto trade-off between accuracy and complexity: as $\lambda$ grows, terms
drop out; the right model sits at the ``knee'', where removing one more term makes the error jump. Choosing
$\lambda$ is model selection (cross-validation, information criteria such as AIC; book Sec.~4.6--4.7).
\end{intuition}

\paragraph{Noise and derivatives.} One rarely measures $\dot{\vb x}$; it must be computed from $\vb x$
\ts{5}{1}{6:58}. With very noisy measurements of $x,y,z$, \emph{total-variation regularized} derivatives
(TVDiff, R.~Chartrand) still allow correct structure identification and an accurate reconstruction of the
attractor \ts{5}{1}{7:18}\ts{5}{1}{7:39}; both steps are numerically robust \ts{5}{1}{7:59}.

\begin{secondary}[Total-variation regularized differentiation]
Find $u\approx\dot f$ as $\min_u\ \alpha\int|u'|\,\mathrm dt+\tfrac12\int\big|\!\int_0^tu-f\big|^2\mathrm dt$:
the integral of the derivative must reproduce the signal, while the total variation of the derivative is
penalized. Unlike finite differences, which amplify noise by $1/\Delta t$, this gives piecewise-smooth
derivatives. Alternatives: Savitzky--Golay filters, local polynomial fits (as in PDE-FIND), or the weak
(integral) form of SINDy, which avoids derivatives altogether.
\end{secondary}

\begin{exercise}[Noise and threshold]
Add Gaussian noise of standard deviation $\eta$ to $X$ in \file{ch05_sindy.py} before differentiating.
For $\eta=0.01,0.1,1$, find the range of $\lambda$ for which the correct structure is found. Replace the
finite differences by a Savitzky--Golay filter (\texttt{scipy.signal.savgol\_filter} with \texttt{deriv=1})
and compare.
\end{exercise}

% ------------------------------------------------------------------
\section{The cylinder wake: rediscovering the mean-field model}\label{sec:si-cylinder}

\begin{casestudy}[Thirty years in a regression]
A benchmark of fluid dynamics: flow past a cylinder, with vortex shedding \ts{5}{1}{8:20}. A low-order model
for it took decades \ts{5}{1}{8:40}:
\begin{enumerate}
  \item Ruelle and Takens (1971), after Landau and Hopf, proposed a route to turbulence through successive
        Hopf bifurcations \ts{5}{1}{9:00}\ts{5}{1}{9:21};
  \item about 15 years later, Zebib, Jackson and others verified numerically that the cylinder wake
        undergoes a Hopf bifurcation at Reynolds number $\approx47$, after which laminar periodic shedding
        sets in \ts{5}{1}{9:42}\ts{5}{1}{10:02};
  \item a paradox: the Hopf normal form \eqref{eq:fu-hopf} has \emph{cubic} terms, but Navier--Stokes has
        only \emph{quadratic} nonlinearity (convection) \ts{5}{1}{10:23}\ts{5}{1}{10:43};
  \item 15 years later again, Noack et al.\ (J.\ Fluid Mech.\ 2003) resolved it with a separation of time
        scales: the \emph{mean-field model}
        \begin{equation}\label{eq:si-meanfield}
          \dot x=\mu x-\omega y+Axz,\qquad \dot y=\omega x+\mu y+Ayz,\qquad \dot z=-\lambda(z-x^2-y^2),
        \end{equation}
        where $x,y$ are the amplitudes of the two shedding modes and $z$ the ``shift mode'' between the
        unstable steady flow and the mean flow. $z$ is fast: it collapses onto the slow parabolic manifold
        $z=x^2+y^2$, on which the dynamics winds out to periodic shedding \ts{5}{1}{11:04}\ts{5}{1}{11:24}\ts{5}{1}{11:46}.
        Only quadratic terms, yet a Hopf bifurcation \ts{5}{1}{12:06}.
\end{enumerate}
SINDy, fed the time series of the first two POD modes and the shift mode from a simulation, identifies a
model with the same structure: only quadratic terms, the slow parabolic manifold, the spiral onto shedding
\ts{5}{1}{12:27}\ts{5}{1}{12:48}\ts{5}{1}{13:10}.
\end{casestudy}

\begin{proposition}[Cubic from quadratic]
On the slow manifold $z=x^2+y^2$, model \eqref{eq:si-meanfield} reduces to the Hopf normal form
$\dot r=\mu r+Ar^3$ (supercritical for $A<0$), $\dot\theta=\omega$.
\end{proposition}
\begin{proof}
With $r^2=x^2+y^2$: $r\dot r=x\dot x+y\dot y=\mu r^2+Ar^2z$, so $\dot r=\mu r+Arz$; substituting $z=r^2$
gives $\dot r=\mu r+Ar^3$. Also $r^2\dot\theta=x\dot y-y\dot x=\omega r^2$. The cubic term is generated by
eliminating the fast variable: adiabatic elimination.
\end{proof}

\paragraph{Kick it off the attractor.} If the data start on the slow manifold, SINDy wrongly identifies
\emph{cubic} terms (the reduced model of the proposition) \ts{5}{1}{13:51}. To get the quadratic model one
must start off the attractor, e.g.\ from the mean flow, and watch the system fall back \ts{5}{1}{14:12}.
General strategy: to identify the dynamics near an attractor, kick the system off it and see how it
returns \ts{5}{1}{14:32}.

\begin{exercise}[Mean-field model]
Simulate \eqref{eq:si-meanfield} with $\mu=0.1$, $\omega=1$, $A=-0.1$, $\lambda=10$ from an initial
condition on the slow manifold, and from one off it ($z$ far from $x^2+y^2$). Run SINDy with polynomials up
to order 3 on both data sets: which one gives the quadratic model, which one the cubic normal form?
\end{exercise}

% ------------------------------------------------------------------
\section{Extensions: parameters, forcing, delays, control}\label{sec:si-extensions}

\paragraph{Parametrized dynamics ``for free''.} Treat a parameter $\mu$ as an extra state with
$\dot\mu=0$ and include it in the library ($\mu,\mu x,\mu x^2,\dots$) \ts{5}{1}{14:53}\ts{5}{1}{15:35}. For the
logistic map, training on 10 values of $\mu$ yields a parametrized model that reproduces the whole
bifurcation diagram \ts{5}{1}{15:14}\ts{5}{1}{15:56}; likewise for the Hopf normal form \ts{5}{1}{15:56}. In the same
way one can add time ($\sin t$, $t^2$, \dots) for non-autonomous systems, or external forcing such as CO$_2$
in a climate model \ts{5}{1}{16:16}\ts{5}{1}{16:36}.

\paragraph{Delay coordinates.} What if only $x$ of the Lorenz system is measured \ts{5}{1}{16:58}\ts{5}{1}{17:18}?
Build the \emph{Hankel matrix} of shifted copies of the time series \ts{5}{1}{17:39}\ts{5}{1}{18:00},
\begin{equation}\label{eq:si-hankel}
  H=\begin{bmatrix}x_1&x_2&\cdots&x_{m-p+1}\\ x_2&x_3&\cdots&x_{m-p+2}\\ \vdots&&&\vdots\\ x_p&x_{p+1}&\cdots&x_m\end{bmatrix}
  =U\Sigma V^*.
\end{equation}
The SVD gives ``eigen time-delay coordinates'': the first columns of $V$ reconstruct an attractor
topologically equivalent to Lorenz's, as Takens' embedding theorem guarantees \ts{5}{1}{18:21}\ts{5}{1}{18:42}.
Running SINDy on the first three columns of $V$ gives a sparse model that reproduces the attractor
\ts{5}{1}{19:02}\ts{5}{1}{19:25}. The same SVD of a Hankel matrix underlies singular spectrum analysis (SSA) in
climate science and the eigensystem realization algorithm (ERA) in system identification; Broomhead and King
(1986) already plotted Lorenz in these coordinates. What is new is a \emph{nonlinear model} on them
\ts{5}{1}{19:45}\ts{5}{1}{20:06}. Chapter~8 (HAVOK) develops this.

\begin{theorem}[Takens, 1981, informal]\label{thm:si-takens}
Let the dynamics evolve on a compact attractor of box dimension $d$ and let $h$ be a generic scalar
measurement. Then for $p>2d$ the delay map $\vb x\mapsto(h(\vb x),h(\vb F_{\tau}\vb x),\dots,h(\vb F_{(p-1)\tau}\vb x))$ is
an embedding: a one-to-one, smooth image of the attractor.
\end{theorem}

\paragraph{SINDy with control.} For $\dot{\vb x}=\vb f(\vb x)+\vb g(\vb u)$, where $\vb u$ may be an
exogenous forcing or a feedback law and may enter nonlinearly ($u^3$, $\sin u$), include $\vb u$ in the
library \ts{5}{1}{20:27}\ts{5}{1}{20:47}\ts{5}{1}{21:08}\ts{5}{1}{21:29}. In the example, Lorenz is trained for $t\in[0,20]$
under a feedback $u=26-x$ plus noise; at $t=20$ the control law is switched. Plain SINDy fails, because it
never knew there was forcing; SINDy with control tracks the new data, because it disentangled the effect of
the control from the dynamics \ts{5}{1}{22:10}\ts{5}{1}{22:31}\ts{5}{1}{22:52}.

\paragraph{The recipe and its open questions} \ts{5}{1}{23:12}\ts{5}{1}{23:33}:
\begin{enumerate}
  \item measure (all the state if possible, otherwise build delay coordinates);
  \item choose a basis (library) in which the dynamics is sparse;
  \item sparse regression for the fewest terms that explain the data.
\end{enumerate}
\begin{itemize}
  \item \emph{Did I measure the right variables?} Delay embeddings give tools to check whether enough
        variables are measured and whether they carry new information \ts{5}{1}{24:35}\ts{5}{1}{24:55}.
  \item \emph{Is the library right?} Harder. Physics guides it (Navier--Stokes means quadratic terms); in
        general, try polynomials, trigonometric functions, rational functions \ts{5}{1}{24:55}\ts{5}{1}{25:17}\ts{5}{1}{25:37}.
  \item The regression is the easy part \ts{5}{1}{25:58}\ts{5}{1}{26:18}.
\end{itemize}

\begin{history}[Ruelle and Takens]
David Ruelle and Floris Takens, \emph{On the nature of turbulence}, Commun.\ Math.\ Phys.\ 20 (1971)
167--192, coined the term \emph{strange attractor} and argued against Landau's (1944) picture of turbulence
as an infinite cascade of Hopf bifurcations: after only a few, a strange attractor generically appears. The
paper was rejected by another journal first and appeared in a journal of which Ruelle was an editor. Takens'
embedding theorem appeared in \emph{Detecting strange attractors in turbulence} (Lecture Notes in Math.\ 898,
1981), the same year as Packard, Crutchfield, Farmer and Shaw's delay reconstruction (PRL 1980).\par
\textit{References:} D.~Ruelle, \emph{Chance and Chaos} (Princeton, 1991); the papers cited.
\end{history}

\begin{exercise}[Parametrized logistic map]
Generate logistic-map data for $\mu\in\{2.5,2.75,\dots,3.95\}$ (10 values), build a library of monomials in
$(x,\mu)$ up to order 4, and fit $x_{k+1}$ (a map: no derivative needed). Use the model to draw the
bifurcation diagram on $[2.5,4]$ and compare with \cref{fig:si-logistic}.
\end{exercise}

\begin{exercise}[Delay embedding]
Using only $x(t)$ from \file{ch03_lorenz.py}, build $H$ with $p=100$ and plot $v_1,v_2,v_3$ (columns of
$V$). How many singular values are significant? Compare the picture with the Lorenz attractor.
\end{exercise}

% ------------------------------------------------------------------
\section{PDE-FIND: discovering partial differential equations}\label{sec:si-pdefind}

The same idea for spatio-temporal data, presented by Samuel Rudy \ts{5}{2}{0:03}. Assume the data solve an
unknown PDE with one time derivative,
\begin{equation}\label{eq:si-pde}
  u_t=N(u,u_x,u_{xx},\dots,Q),
\end{equation}
where $Q$ is extra information (a potential, a velocity field), and that $N$ is sparse in a guessed basis,
e.g.\ polynomials in $u$ and its derivatives \ts{5}{2}{0:24}\ts{5}{2}{0:45}\ts{5}{2}{1:06}.

\paragraph{Steps} \ts{5}{2}{1:48}--\ts{5}{2}{3:35}:
\begin{enumerate}
  \item compute $u_t$ and the spatial derivatives from the discretized data: finite differences if the
        data are clean; with noise, a smoothing method (local polynomial interpolation works, but choosing
        the number of points and the degree is tricky) \ts{5}{2}{2:08};
  \item stack every space-time point as one row: $U_t\approx\Theta(U,Q)\,\boldsymbol\xi$, with columns
        such as $1,u,u_x,u^2,uu_x,\dots$ (nothing forbids $\cos u$ if physics suggests it)
        \ts{5}{2}{2:30}\ts{5}{2}{2:50}\ts{5}{2}{3:10};
  \item solve for a sparse $\boldsymbol\xi$.
\end{enumerate}
Least squares always gives dense solutions (a 100-term PDE whose terms are mostly noise), and the
columns of $\Theta$ are correlated, so the matrix is ill-conditioned \ts{5}{2}{3:35}\ts{5}{2}{3:56}\ts{5}{2}{4:17}.
Among LASSO, elastic net, greedy methods and thresholding, Rudy chose \emph{sequential threshold ridge
regression} (STRidge): ridge estimate, cut coefficients below a threshold, repeat on the survivors
\ts{5}{2}{4:17}\ts{5}{2}{4:37}. For Navier--Stokes the sparse fit keeps the 4 correct terms, least squares
keeps all 60 \ts{5}{2}{4:58}\ts{5}{2}{5:19}; the coefficients are within 1\% \ts{5}{2}{6:24}.

\codefile{PDE-FIND on Burgers' equation}{../code/ch05_pdefind.py}

The script generates Burgers' equation $u_t=-uu_x+0.1u_{xx}$ spectrally, computes all derivatives with
finite differences, and from a library of 12 terms finds $u_t=0.1016\,u_{xx}-1.0008\,uu_x$; with a random
1\% of the rows, $u_t=0.1016\,u_{xx}-1.0030\,uu_x$.

\paragraph{Sub-sampling.} The number of rows is the number of data points, about 13 million for a 2D
Navier--Stokes data set in time; using about 1\% of them (dropping rows) gives the same result
\ts{5}{2}{6:45}\ts{5}{2}{7:06}.

\paragraph{Which PDE? The soliton ambiguity.} A function may solve many PDEs \ts{5}{2}{7:30}. A single KdV
soliton, $u=\frac c2\operatorname{sech}^2\!\big(\frac{\sqrt c}{2}(x-ct)\big)$, also solves the advection
equation $u_t=-cu_x$, and PDE-FIND returns the simpler one, advection \ts{5}{2}{7:50}. Two solitons of
different speeds in one solution, or two separate solutions with solitons of different amplitudes, force a
nonlinear PDE, and PDE-FIND finds KdV \ts{5}{2}{8:10}\ts{5}{2}{8:31}\ts{5}{2}{8:52}. The lesson: the data must
be rich enough to exclude the alternatives, as in ``kick it off the attractor''.

\paragraph{Gallery.} KdV, Burgers, nonlinear Schr\"odinger, linear Schr\"odinger in a harmonic potential,
Kuramoto--Sivashinsky, reaction--diffusion, vorticity equation of Navier--Stokes: all recovered
\ts{5}{2}{9:13}\ts{5}{2}{9:33}. With noise, Kuramoto--Sivashinsky keeps the right sparsity pattern but its
coefficients are badly underestimated (52\% error) \ts{5}{2}{9:33}. Noise complicates everything through the
derivatives; next steps are real data and latent variables \ts{5}{2}{9:55}\ts{5}{2}{10:15}.

\begin{exercise}[Soliton ambiguity]
Generate one KdV soliton and run \file{ch05_pdefind.py} with a library containing $u_{xxx}$ and $uu_x$.
Explain why the regression cannot distinguish $u_t=-cu_x$ from $u_t=-6uu_x-u_{xxx}$, and verify that both
are satisfied (choose the soliton normalization accordingly).
\end{exercise}

\begin{keyformulas}
\begin{itemize}
  \item SINDy: $\dot X\approx\Theta(X)\Xi$, $\Xi$ sparse (STLSQ or LASSO), one regression per variable.
  \item Library extensions: parameters ($\dot\mu=0$), time, forcing $\vb u$ (SINDYc), delay coordinates
        from the SVD of the Hankel matrix.
  \item Mean-field model: quadratic terms + fast variable $\Rightarrow$ cubic Hopf on $z=x^2+y^2$.
  \item PDE-FIND: $U_t\approx\Theta(U,U_x,U_{xx},\dots)\boldsymbol\xi$, STRidge; derivatives are the hard part.
\end{itemize}
\end{keyformulas}

\end{paracol}
